\documentclass[12pt]{article}
\usepackage[utf8]{inputenc}
\usepackage[T1]{fontenc}
\usepackage{hyperref}
\title{Kuku-Mu'inh Language Resources for K-12 Curriculum Development}
\author{Compiled by Claude Sonnet 5.5 (claude-sonnet-5-5) for Andr\'as Kornai}
\date{}
\begin{document}
\maketitle

\section{Introduction}
Kuku-Mu'inh (ISO 639-3 xmp; Glottocode kuku1282) is a traditional language of Cape York Peninsula, Queensland, Australia. Search results place it within the Wik/Kugu Paman language complex. Sources give it as a variety of Kugu Nganhcara, also called Kugu Muminh or Wik Muinh, and say the country lies within Cook Shire. The language may no longer have native speakers. The sources are mostly Wikipedia and mirrors, so the naming and dialect boundaries are not settled, and this report should be read with that in mind. Related varieties carry other codes (xmh Kuku-Muminh, ugb Kuku-Ugbanh), and I could not separate resources for them cleanly from those for xmp.

\section{Primary Educational Resources}
None found for Kuku-Mu'inh itself.

\section{Linguistic and Reference Materials}
No dictionary or grammar of Kuku-Mu'inh turned up. The following describe the wider complex:
\begin{itemize}
\item Kugu Nganhcara language (Wikipedia): \url{https://en.wikipedia.org/wiki/Kugu_Nganhcara_language}
\item Kugu Muminh (Wikipedia mirror): \url{https://troywiki.duckdns.org/content/wikipedia_en_all_maxi_2024-01/A/Kugu-Muminh}
\item Edward River, Queensland (locality page naming Kugu Mu'inh as a traditional language of the area): \url{https://en.wikipedia.org/wiki/Edward_River,_Queensland}
\end{itemize}

\section{Religious and Community Resources}
None identified.

\section{Digital and Audio Resources}
\begin{itemize}
\item Endangered Languages Project entries surfaced for the area; I did not verify which variety they document: \url{https://endangeredlanguages.com/lang/3446/samples/11510}
\item Lund University Cape York field-site page, relating to documentation of highly endangered Paman languages: \url{https://projekt.ht.lu.se/lacola/field-sites/cape-york-australia/}
\end{itemize}

\section{Curriculum Development Context}
Revitalisation work would be a heritage and revival effort led by the descendant community. This report found no school programme, community language centre or archive of materials for Kuku-Mu'inh specifically. The first step would be to find the existing documentation, probably through the collaborative documentation project on Cape York languages, and to contact local Wik and Kugu community bodies. The searches did not identify these contacts, so I name none. Because Kugu Muminh and Kugu Uwanh are described as the dialects still spoken, nearby resources may be adaptable, subject to community consent.

\end{document}